% paper/sections/conclusion.tex

\section{Conclusion}
\label{sec:conclusion}

We presented Rev-GNN-LSTM, the first end-to-end learned policy for joint seed
selection and personalized pricing in social networks.
By coupling a GraphSAGE encoder with an LSTM episode-history module and training
via imitation learning followed by REINFORCE, the model captures the joint
dependency between targeting and pricing decisions that hand-crafted baselines
ignore.
On Forest-Fire $n{=}1000$ the model achieves $462.6$ revenue versus $460.0$
for the best hand-crafted baseline ($p{<}0.0001$, $20/20$ seeds), %SOURCE: CLAUDE.md
and generalizes zero-shot to unseen network sizes and topologies.

We further introduced a budget-constrained formulation — production cost and
replenishable capital — that transforms the problem into a sequential knapsack
with network externalities.
A family of calibrated dynamic-programming baselines (Cal-DP v2, v3, composite)
dominates the prior standard at every budget level tested.
The learned policy excels in the capital-binding regime ($k \le 5$, up to $41\times$
vs.\ Greedy+Budget at $k{=}1$), while Cal-DP leads at mid-$k$ on dense real
networks — a finding we report as a honest regime-aware conclusion.

\paragraph{Future work.}
Direct extensions include (i) training the budget-constrained policy on denser
graphs closer to real social networks; (ii) replacing the Rayleigh valuation model
with a learned model calibrated from observed purchase histories;
(iii) extending to multi-item assortment settings where each buyer chooses from
a menu; and (iv) applying the framework to decentralized recommendations where
the ``seller'' is a platform and seeds are content creators.
The formalization of revenue maximization as an RL problem opens a path toward
data-driven market design at social-network scale.
